% SEGMENT OLP-0403-S01
% Part: many-valued-logic
% Chapter: sequent-calculus
% Section: introduction

\documentclass[../../../include/open-logic-section]{subfiles}

\begin{document}

\olfileid{mvl}{seq}{int}

\olsection{அறிமுகம்}

செவ்வியல் தருக்கத்திற்கான தொடரணிக் கணியம் திறம்பட்ட, எளிய
!!{derivation} முறைமை. ஒரு பல்மதிப்புத் தருக்கம் முடிவுறு எண்ணிக்கையிலான
மெய்மதிப்புகளைக் கொண்ட ஓர் அணியால், அதாவது முடிவுறு~$V$ ஆல்,
வரையறுக்கப்பட்டிருந்தால், அதற்கு ஒரு தொடரணிக் கணியத்தை வழங்க முடியும்.
செவ்வியல் நேர்வில் தொடரணிகளின் பொருளையும் அனுமான விதிகளின் வடிவத்தையும்
கருதுவதிலிருந்து இதை எவ்வாறு செய்வது என்ற எண்ணம் கிடைக்கிறது.

ஒரு தொடரணி
% SOURCE CORRECTION TA-MVL-021: use A_m as the final antecedent to match the interpreted conjunction.
\begin{align*}
    !A_1, \dots, !A_m & \Sequent !B_1, \dots, !B_n
\intertext{பின்வரும் !!{formula} ஆகப் பொருள்கொள்ளப்படலாம் என்பதை இப்போது
நினைவுகூர்க:}
    (!A_1 \land \cdots \land !A_m) & \lif (!B_1 \lor \cdots \lor
!B_n)
\end{align*}
வேறு சொற்களில், ஒரு தொடரணி~$\Gamma \Sequent \Delta$ ஐ
!!^a{valuation}~$\pAssign{v}$ \emph{நிறைவுறுத்துவது}, அப்போதும் அப்போது
மட்டுமே, ஏதோ~$!A \in \Gamma$ க்கு $\pValue{v}(!A) = \False$ அல்லது
ஏதோ~$!A \in \Delta$ க்கு $\pValue{v}(!A) = \True$. இந்தப்
பொருள்கோளின்படி, தொடக்கத் தொடரணிகள்~$!A \Sequent !A$ எப்போதும்
நிறைவுறுத்தப்படுகின்றன; ஏனெனில் $\pValue v(!A) = \True$ அல்லது
% SOURCE CORRECTION TA-MVL-022: restore the valuation argument in the false-value alternative.
$\pValue v(!A) = \False$.

பக்க வாய்பாடுகள்~$\Gamma$, $\Delta$ விடப்பட்ட நிலையில்,~$\Log{LK}$ இல்
நிபந்தனைவாக்கியத்திற்கான அனுமான விதிகள் இவை:

\begin{defish}
    \Axiom$ \fCenter !A$
    \Axiom$ !B \fCenter $
    \RightLabel{\LeftR{\lif}}
    \BinaryInf$ !A \lif !B \fCenter $
    \DisplayProof
    \hfill
    \Axiom$ !A \fCenter !B$
    \RightLabel{\RightR{\lif}}
    \UnaryInf$ \fCenter !A \lif !B $
    \DisplayProof
\end{defish}

மேலுள்ள ஒரு தொடரணியின் பொருண்மைப் பொருள்கோளைப் பயன்படுத்தினால்,
$\pValue v(!A) = \True$, $\pValue v(!B) = \False$ எனில்,
$\pValue v(!A \lif !B) = \False$ என்று $\LeftR{\lif}$ விதி கூறுவதாகப்
படிக்கலாம். இதேபோல், $\pValue v(!A) = \False$ அல்லது
$\pValue v(!B) = \True$ எனில், $\pValue v(!A \lif !B) = \True$ என்று
$\RightR{\lif}$ விதி கூறுகிறது. உண்மையில், இந்த நிபந்தனைக் கூற்றுகள்
இருவழிக் கூற்றுகளே. $\LeftR{\land}$, $\RightR{\lor}$ விதிகளின் செந்தர
வடிவங்களில், அவற்றிற்கு ஒத்த நிபந்தனைக் கூற்றுகள் இருவழிக் கூற்றுகளாக
இருக்காது. ஆனால் அவை இருவழிக் கூற்றுகளாக அமையும் பின்வரும் மாற்று வடிவங்கள்
இந்த விதிகளுக்கு உள்ளன:

\begin{defish}
    \Axiom$!A, !B, \Gamma \fCenter \Delta$
    \RightLabel{\LeftR{\land}}
    \UnaryInf$!A \land !B, \Gamma \fCenter \Delta$
    \DisplayProof
    \hfill
    \Axiom$ \Gamma \fCenter \Delta, !A, !B$
    \RightLabel{\RightR{\lor}}
    \UnaryInf$ \Gamma \fCenter \Delta, !A \lor !B$
    \DisplayProof
\end{defish}

இந்த அடிப்படை எண்ணத்தை ஓர்~$n$-மதிப்புத் தருக்கத்திற்குப் பயன்படுத்தினால்,
இரண்டு இடங்களுக்குப் பதிலாக, ஒவ்வொரு மெய்மதிப்புக்கும் ஒன்றாக, $n$ இடங்களைக்
கொண்ட ஒரு தொடரணிக் கணியம் கிடைக்கிறது. $V = \{\False, \Undef, \True\}$
என்ற மும்மதிப்புத் தருக்கத்தில், ஒரு தொடரணி
$\Gamma \mid \Pi \mid \Delta$ என்ற கோவை. அதை !!a{valuation}~$\pAssign v$
நிறைவுறுத்துவது, அப்போதும் அப்போது மட்டுமே, ஏதோ~$!A \in \Gamma$ க்கு
$\pValue{v}(!A) = \False$ அல்லது ஏதோ~$!A \in \Delta$ க்கு
$\pValue{v}(!A) = \True$ அல்லது ஏதோ~$!A \in \Pi$ க்கு
$\pValue{v}(!A) = \Undef$. எனவே தொடக்கத் தொடரணிகள்
$!A \mid !A \mid !A$ எப்போதும் நிறைவுறுத்தப்படுகின்றன.

\end{document}
